%%%%%%%%%%%%%%%%%%%%%%%%%%%%%%%%%%%%%%%%%
% Classicthesis-Styled CV
% LaTeX Template
% Version 1.0 (22/2/13)
%
% This template has been downloaded from:
% http://www.LaTeXTemplates.com
%
% Original author:
% Alessandro Plasmati
%
% License:
% CC BY-NC-SA 3.0 (http://creativecommons.org/licenses/by-nc-sa/3.0/)
%
%%%%%%%%%%%%%%%%%%%%%%%%%%%%%%%%%%%%%%%%%

%----------------------------------------------------------------------------------------
%	PACKAGES AND OTHER DOCUMENT CONFIGURATIONS
%----------------------------------------------------------------------------------------

\documentclass{scrartcl}

\reversemarginpar % Move the margin to the left of the page 

\newcommand{\MarginText}[1]{\marginpar{\raggedleft\itshape\small#1}} % New command defining the margin text style

\usepackage[nochapters]{classicthesis} % Use the classicthesis style for the style of the document
\usepackage[LabelsAligned]{currvita} % Use the currvita style for the layout of the document
\usepackage{graphicx}



\renewcommand{\cvheadingfont}{\LARGE\color{Maroon}} % Font color of your name at the top

\newcommand{\cemph}{\it\color{Maroon}}

\usepackage{hyperref} % Required for adding links	and customizing them
\hypersetup{colorlinks, breaklinks, urlcolor=Maroon, linkcolor=Maroon} % Set link colors

\newlength{\datebox}\settowidth{\datebox}{Spring 2011} % Set the width of the date box in each block

\newcommand{\NewEntry}[3]{\noindent\hangindent=2em\hangafter=0 \parbox{\datebox}{\small \textit{#1}}\hspace{1.5em} #2 #3 % Define a command for each new block - change spacing and font sizes here: #1 is the left margin, #2 is the italic date field and #3 is the position/employer/location field
\vspace{0.5em}} % Add some white space after each new entry

\newcommand{\Description}[1]{\hangindent=2em\hangafter=0\noindent\raggedright\footnotesize{#1}\par\normalsize\vspace{1em}} % Define a command for descriptions of each entry - change spacing and font sizes here
\newcommand{\MoreDescription}[1]{\hangindent=2em\hangafter=0\noindent\raggedright\scriptsize{#1}\par\normalsize\vspace{1em}} % Define a command for descriptions of each entry - change spacing and font sizes here

%----------------------------------------------------------------------------------------

\begin{document}

\thispagestyle{empty} % Stop the page count at the bottom of the first page

%----------------------------------------------------------------------------------------
%	NAME AND CONTACT INFORMATION SECTION
%----------------------------------------------------------------------------------------

\begin{cv}{Tim {\Large CD} Lucas}\vspace{1.5em} % Your name

\noindent{\color{Maroon}\spacedlowsmallcaps{Personal Information}}\vspace{0.5em} % Personal information heading

%\NewEntry{}{\textit{Born in Canada,}}{05 September 1987} % Birthplace and date

\NewEntry{email}{\includegraphics[height=8pt]{Ar_Icon_Contact} \href{mailto:timcdlucas@gmail.com}{timcdlucas{\small @}gmail.com}} % Email address

\NewEntry{twitter}{\includegraphics[height=8pt]{Twitter_logo_blue-small.png} \href{http://www.twitter.com/timcdlucas}{{\small @}timcdlucas}  \href{http://www.twitter.com/statsforbios}{{\small @}statsforbios}} % Personal website

\NewEntry{website}{\href{http://www.ucl.ac.uk/~ucbptcl}{www.ucl.ac.uk/{\raise.17ex\hbox{$\scriptstyle\mathtt{\sim}$}}ucbptcl}} % Personal website

\NewEntry{github}{\href{http://www.github.com/timcdlucas}{www.github.com/timcdlucas}} % Personal website

\NewEntry{scholar}{\href{http://scholar.google.co.uk/citations?user=WfpSfMAAAAAJ&hl=en&oi=ao}{Google scholar}} % Personal website



\NewEntry{phone}{07415 863 536} % Phone number(s)

\vspace{1em} % Extra white space between the personal information section and goal

\hspace{-0.4cm}{\color{Maroon}\spacedlowsmallcaps{Present Appointment}}\vspace{1em}

\NewEntry{2016--present}{CBER Programmer}

\Description{\MarginText{Research Programmer}I am currently the stuff programmer for the Centre of Biodiversity and Environment Research at UCL. I work on two main projects. I am working with the \href{http://www.madingleymodel.org/}{Madingley Model}---an ecological model of all life, written in C\#. My primary task is to get this model running on the {\cemph high performance cluster} at UCL. Secondly, I am translating code from {\cemph Mathematica to R} that performs statistical analyses for measurements of 3D objects used in {\cemph paleontological research}. I also provide technical support for the rest of the department.}

%------------------------------------------------

\vspace{1em} % Extra space between major sections


%----------------------------------------------------------------------------------------
%	EDUCATION
%----------------------------------------------------------------------------------------

{\color{Maroon}\spacedlowsmallcaps{Education}}\vspace{1em}

\NewEntry{2012--present}{University College London, {\small CoMPLEX}}

\Description{\MarginText{PhD}\textit{Social structure and network epidemiology in bat zoonoses}\newline
Description: I am using {\cemph complex networks} to study the epidemiology of bat-borne diseases. As bats carry a number of important {\cemph zoonotic diseases}, understanding the spread of these diseases within the bat population and how this affects spillover to humans and livestock is increasingly important. The unusually social nature of bat populations will strongly affect how diseases spread.\newline
Supervisors: Prof.~Kate \textsc{Jones} \& Dr~Hilde \textsc{Wilkinson-Herbots}}

%------------------------------------------------

\NewEntry{2011--2012}{University College London, {\small CoMPLEX}}

\Description{\MarginText{MRes}\textit{Modelling Biological Complexity} $\cdotp$ {\bf Merit}\newline
Description: This was part of a combined MRes/PhD program. It was an interdisciplinary course applying {\cemph quantitative methods} to the life sciences. }

%------------------------------------------------

\NewEntry{2006--2010}{University of Sheffield, {\small Animal \& Plant Sciences}}

\Description{\MarginText{MBioSci}\textit{Zoology} $\cdotp$ {\bf First}\newline
Description: For my final project I used {\cemph wavelet} analysis to study multi-annual {\cemph cycles in malaria} incidence in Thailand.}

%------------------------------------------------
\vspace{1em} % Extra space between major sections


%----------------------------------------------------------------------------------------
%	WORK EXPERIENCE
%----------------------------------------------------------------------------------------

\noindent{\color{Maroon}\spacedlowsmallcaps{Research Experience}}\vspace{1em}




%------------------------------------------------

\NewEntry{Autumn 2014}{Zo\"{o}n: An R package for reproducible SDMs}

\Description{\MarginText{Internship}I wrote the first version of an R package for {\cemph reproducible} species distribution modelling. The package uses an online repository of user submitted `modules' to allow the software to keep up with this fast moving field and allow analyses to be completely reproducible. [\href{http://www.github.com/zoonproject/zoon}{Github}]}

\NewEntry{Summer 2012}{Estimating abundances using acoustic data}

\Description{\MarginText{Summer Project}I adapted `{\cemph ideal gas}' models to acoustic data. I applied the model using R to a pan-European bat survey. We have worked on this project further, validating results with simulations, and the work is now published.}

%------------------------------------------------

\NewEntry{May 2012}{Pair approximations in spatial biology}

\Description{\MarginText{Case Presentation}I compared a number of moment closures for a {\cemph pair-approximation} model of tree population growth to lattice simulations written in Mathematica. [\href{http://files.figshare.com/92192/Pair_approximations_in_spatial_biology_Tim_Lucas_2012.pdf}{pdf}]}


%------------------------------------------------

\NewEntry{Dec. 2011}{Gaussian processes for bat identification}

\Description{\MarginText{Case Presentation}I applied a novel {\cemph machine learning} method to a library of {\cemph bat calls} in Matlab. I compared the effectiveness of this method to standard machine learning methods applied in R. [\href{http://www.ucl.ac.uk/~ucbptcl/bat-gaussian-process-learning-tim-lucas-2011.pdf}{pdf}]}

%------------------------------------------------

\NewEntry{August 2011}{Smithsonian Tropical Research Institute}

\Description{\MarginText{Volunteer Fieldwork}Two months fieldwork in Panam\'{a} on two projects: studying \emph{Anolis} dewlap evolution and studying gut length plasticity in Red-eyed tree frogs. }

%------------------------------------------------

\NewEntry{May 2011}{Chilo\'{e} Silvestre, Chil\'{e}}

\Description{\MarginText{Volunteer Fieldwork}I spent two weeks trapping Darwin's foxes in Chil\'{e} to collect samples for geophylogenetics.}

%------------------------------------------------

\NewEntry{August 2010}{University of Sheffield}

\Description{\MarginText{Summer Internship}I studied the evolutionary response of plant communities to climate change with Dr Raj Whitlock. I collected, propagated and analysed plants collected from the field. }

%------------------------------------------------

\NewEntry{August 2009}{University of York, {\small YCCSA}}

\Description{\MarginText{{\footnotesize TRANSIT} Internship}I studied collective foraging behaviour by programming a {\cemph complex 3{\footnotesize D} foraging} {\cemph simulation} in Java and running simulations on a cluster at the York Centre for Complex Systems Analysis.}

%------------------------------------------------

\vspace{2em} % Extra space between major sections



%----------------------------------------------------------------------------------------
%	PUBLICATIONS
%----------------------------------------------------------------------------------------
{\color{Maroon}\spacedlowsmallcaps{Publications}}\vspace{1em}


\Description{\MarginText{2016}{\bf Lucas TCD}, Wilkinson-Herbots HM, \& Jones KE. \emph{A comparative and computational study of population structure and pathogen richness in bats.} In prep.}

\Description{\MarginText{2015}{\bf Lucas TCD}\textsuperscript{$\ast$}, Moorcroft EA\textsuperscript{$\ast$}, Freeman R, Rowcliffe MJ \& Jones KE. (2015) \emph{A generalised random encounter model for estimating animal density with remote sensor data.} Methods in Ecology and Evolution. doi: 10.1111/2041-210X.12346 [\href{http://onlinelibrary.wiley.com/doi/10.1111/2041-210X.12346/epdf}{pdf}]}

\Description{\MarginText{2013}Walters CL, Collen A, {\bf Lucas TCD}, Mroz K, Sayer CA and Jones KE. (2013) Challenges of Using Bioacoustics to Globally Monitor Bats. in \emph{Bat Evolution, Ecology, and Conservation.} Springer New York. 479-499.}

\Description{\MarginText{} \scriptsize{ $\ast$ Co-first authors. } }

%------------------------------------------------

\vspace{1em} % Extra space between major sections


%----------------------------------------------------------------------------------------
%	Software
%----------------------------------------------------------------------------------------

{\color{Maroon}\spacedlowsmallcaps{Software}}\vspace{1em}


\Description{\MarginText{On CRAN}{\bf Tim Lucas}, Nick Golding, Tom August, Greg McInerny, Emiel van Loon (2015) \emph{Zo\"{o}n: Reproducible, Accessible \& Shareable Species Distribution Modelling}. \url{www.github.com/zoonproject/zoon}}

\Description{{\bf Tim Lucas} (2015) \emph{palettetown: Use Pokemon Inspired Colour Palettes} \url{www.github.com/timcdlucas/palettetown}}



%------------------------------------------------

\vspace{2em} % Extra space between major sections

%----------------------------------------------------------------------------------------
%	OTHER INFORMATION
%----------------------------------------------------------------------------------------

{\color{Maroon}\spacedlowsmallcaps{Conferences}}\vspace{1em}

\Description{\MarginText{Meetings 2015}\emph{The Zo\"{o}n Project: Reproducible, Remixable and Shareable Species Distribution Modelling with R.}}\vspace{-1em}
\MoreDescription{Presentation at \href{http://www.britishecologicalsociety.org/events/current_future_meetings/past-bes-annual-meetings/2015-annual-meeting/}{BES Annual Meeting} by T. August, N. Golding, \textbf{T. Lucas}, D. Gavaghan, N. Isaac, B. O'Hara, E. van Loon \& G. McInerny  }

\vspace{-0.5em}


\Description{\emph{Simple, Shareable and Reproducible Species Distribution Modelling with the Zo\"{o}n R package.}}\vspace{-1em}
\MoreDescription{Poster at \href{http://www.britishecologicalsociety.org/events/current_future_meetings/past-bes-annual-meetings/2015-annual-meeting/}{BES Annual Meeting} by N. Golding, \textbf{T. Lucas}, T. August, D. Gavaghan, N. Isaac, B. O'Hara, E. van Loon \& Greg McInerny  }

\vspace{-0.5em}

\Description{\emph{Comparative and computational studies of pathogen richness in bats.}}\vspace{-1em}
\MoreDescription{Presentation at Research in Progress, \href{https://rstmh.org/events/research-progress-2015}{RSTMH} by \textbf{T. Lucas}, H. Wilkinson-Herbots \& K. Jones.}

\vspace{-0.5em}

\Description{\emph{A comparative and computational study of population structure and pathogen richness in bats.}}\vspace{-1em}
\MoreDescription{Presentation at \href{http://www.epidemics.elsevier.com/}{Epidemics5} conference by \textbf{T. Lucas}, H. Wilkinson-Herbots \& K. Jones.}

\vspace{-0.5em}



\Description{\emph{An ideal gas model for estimating absolute abundances from bat detector data.}}\vspace{-1em}
\MoreDescription{Presentation at the National Bat Conference. [\href{http://www.slideshare.net/timcdlucas/tim-lucasnbc}{slides}]}

\vspace{-0.5em}

\Description{\emph{Pathogen diversity and bat population structure.}}\vspace{-1em}
\MoreDescription{Poster at British Parasitological Society Autumn Meeting.}


\vspace{-0.5em}

\Description{\emph{Estimating abundance from camera traps and acoustic sensors.}}\vspace{-1em}
\MoreDescription{Presentation at {\footnotesize CEH}, Wallingford seminar series.}


\vspace{-0.5em}

\Description{\MarginText{2014}Presentation at \href{http://id2-ox.co.uk/}{id2oxford} conference. [\href{http://www.slideshare.net/timcdlucas/tim-lucasid2ox}{slides}]}
\vspace{-0.5em}

\Description{Poster at the {\footnotesize CoMPLEX} conference. [\href{http://figshare.com/articles/Estimating_abundance_from_camera_traps_and_acoustic_sensors/1321269}{pdf}]}
\vspace{-0.5em}

\Description{\MarginText{2013}Presentation at BritBats 2 [\href{http://www.slideshare.net/timcdlucas/tim-lucasbritbats}{slides}].}

\vspace{-0.5em} % Negative vertical space to counteract the vertical space between every \Description command

\Description{Invited attendance at \href{http://ecoviz.wordpress.com/1-joining-things-up-and-filling-the-gaps/}{eco{\footnotesize VIZ}} Tansley workshop.}

\vspace{-0.5em} % Negative vertical space to counteract the vertical space between every \Description command

\Description{Poster at the {\footnotesize CoMPLEX} conference and id2 conference. [\href{http://figshare.com/articles/The_Phylogeography_of_Henipah_Virus/1046697}{pdf}]}

\Description{\MarginText{Teaching}2015\ \ $\cdotp$\ \ Demonstrator for reproducible species distribution modelling workshop run by Quantitative Ecology special interest group at BES.}
\vspace{-0.5em} % Negative vertical space to counteract the vertical space between every \Description command

\Description{2013--2014\ \ $\cdotp$\ \ Online tutor for \href{http://sysmic.ac.uk/home.html}{{\footnotesize SysMIC}}, a course for teaching quantitative skills to biologists.}

\Description{\MarginText{Peer Review}\emph{Journals Reviewed for:}}

\vspace{-0.5em} % Negative vertical space to counteract the vertical space between every \Description command

\Description{\ \ $\cdotp$\ \ Methods in Ecology and Evolution, National Academy Science Letter}



%------------------------------------------------



%----------------------------------------------------------------------------------------
%	COMPUTER SKILLS
%----------------------------------------------------------------------------------------

{\color{Maroon}\spacedlowsmallcaps{Computational Skills}}\vspace{1em}

\Description{\MarginText{Languages}R (eight years), Python, Matlab, Mathematica, Java, SQL.}

\Description{\MarginText{OS}Comfortable with Windows, Mac or Linux.}

\Description{\MarginText{Other}Experience in Git/Github, unit testing, LaTeX, web design, markdown, R package development, shell/ssh and high performance computing.}

\end{cv}

\end{document}
